%% ========================================================================
%% Chapter 3: Dasar Input/Output (I/O)
%% ========================================================================

\chapter{Dasar Input/Output (I/O)}
\label{ch:basic-io}

Bab ini membahas mekanisme input/output di Linux, termasuk file descriptor,
pengalihan I/O, pipa, dan penggunaan operator untuk manipulasi aliran data.

%% ------------------------------------------------------------------------
\section{File Descriptor di Linux}
\label{sec:file-descriptors}

% [Konten akan diisi di sini]

%% ------------------------------------------------------------------------
\section{Pengalihan I/O Standar}
\label{sec:io-redirection}

\subsection{Standard Input (stdin)}
\subsection{Standard Output (stdout)}
\subsection{Standard Error (stderr)}

% [Konten akan diisi di sini]

%% ------------------------------------------------------------------------
\section{Pipa dan Operator |}
\label{sec:pipes}

% [Konten akan diisi di sini]

%% ------------------------------------------------------------------------
\section{Pengalihan Output ke File}
\label{sec:output-redirection}

\subsection{Operator > (Overwrite)}
\subsection{Operator >> (Append)}

% [Konten akan diisi di sini]

%% ------------------------------------------------------------------------
\section{Penggunaan tee untuk Output Ganda}
\label{sec:tee-command}

% [Konten akan diisi di sini]

%% ------------------------------------------------------------------------
\section{Rangkuman}
\section{Latihan}
\section{Referensi Tambahan}
